\documentclass[10pt]{article}
\usepackage[margin=20mm]{geometry}
\usepackage{amsmath,amssymb,mathtools,fontspec,unicode-math,xurl,hyperref}
\IfFontExistsTF{DejaVu Serif}{\setmainfont{DejaVu Serif}}{\setmainfont{Latin Modern Roman}}
\setmathfont{latinmodern-math.otf}
\setlength{\parindent}{0pt}\setlength{\parskip}{6pt}\sloppy
\title{Technical Verification of
The Connected Components of Affine Deligne–Lusztig Varieties}
\author{AI-assisted technical review}
\date{}
\begin{document}
\maketitle
Abstract: This claim-focused initial review establishes that the component-quotient assertion of Lemma 3.2 is false in its stated generality. An explicit diamond quotient has a connected underlying space although the action on source components is not transitive. The lemma is the stated justification of Proposition 3.12, which is used in the central implication of Theorem 6.1. This supports a qualified Major finding about a core proof dependency, without a counterexample to the actual shtuka assertions or the main component formula. A separate converse argument in Proposition 6.6 remains a Needs Evidence item. Scientific adjudication and independent acceptance are pending.\par
\section*{Paper identity and inspected scope}
The reviewed paper is Ian Gleason, Dong Gyu Lim and Yujie Xu, The connected components of affine Deligne–Lusztig varieties, DOI https://doi.org/10.1007/s00222-025-01386-1 [GLX]. The canonical source is the acquired publisher version of record, 57 PDF pages, printed pages 805–861. It was published online on 12 November 2025 and assigned to volume 243 (2026). Throughout this report, unqualified page numbers are printed pages; physical PDF page \(j\) is printed page \(804+j\).\par
The review checks the main theorem’s hypotheses and reduction structure, the component-quotient lemma and its exact uses, the ad-isomorphism reduction, coarsening of parahoric level, and the final component-stabilizer argument. The long geometric constructions in Sections 4–5 and the integral-model applications are not independently rederived. Nineteen source pages supporting these checks were inspected as rendered images; the coverage ledger distinguishes this from other extracted-text reading.\par
The acquired arXiv source is v3 of arXiv:2208.07195 (10 November 2025); it is supporting author material, not a replacement for the version of record. No separately applicable supplement or computational code was identified in the bounded source search. Code absence is not a mathematical defect. A live title/DOI search did not identify a formal correction; this is a bounded search result, not proof that none exists.\par
\section*{Definitions, hypotheses and dependencies}
Equation reference: eq:published. Immediately before Lemma 3.2 (p. 820), the paper defines \(\pi_0(F)\) as the connected components of the underlying topological space \(|F|\), equipped with the quotient topology. Lemma 3.2 assumes that \(K\) is locally profinite, \(F,G\) are small \(v\)-sheaves, and \(G=F/\underline K\). It asserts \[
 |G|=|F|/K,\qquad \pi_0(G)=\pi_0(F)/K.\] It does not assume a compact group, a proper action, quasiseparatedness, or local spatiality.\par
The main geometric applications have additional structure. In particular, actual infinite-level shtuka spaces and their period maps are not arbitrary small \(v\)-sheaves, and the Iwahori and Galois groups occurring in several quotient steps are compact. The counterexample below must be evaluated against the written lemma, while these additional structures limit what it says about those applications.\par
The relevant core dependency is explicit. Proposition 3.12 (p. 828) invokes Lemma 3.2 to deduce transitivity of \(G(\mathbf Q_p)\) on shtuka components from connectedness of the admissible locus. Proposition 6.11 (pp. 854–855) invokes Proposition 3.12 to identify that component set with a homogeneous space \(G(\mathbf Q_p)/H_t\). Lemmas 6.12–6.14 then use this description in the central \((4)\Rightarrow(5)\) implication of Theorem 6.1. This is a concrete dependency in the proof.\par
\section*{Confirmed substantive finding}
Problem 1 (The component-quotient assertion in Lemma 3.2).  The second equality in [eq:published] (\#eq:published) fails for an actual free action of a locally profinite group on an affinoid perfectoid space, with a diamond as quotient.\par
\subsection*{A counterexample satisfying the stated assumptions}
Equation reference: eq:example. Set \(p=2\). Let \(C\) be an algebraically closed complete perfectoid field of characteristic \(2\), for example the completion of an algebraic closure of \(\mathbf F_2((t))\), and let \(C^+=\mathcal O_C\). Put \[
\begin{split}
 S&=\mathbf Z_2,\qquad K=\mathbf Z_{\mathrm{discrete}},\\
 F&=\underline S\times\operatorname{Spa}(C,\mathcal O_C)
   =\operatorname{Spa}\!\left(C^0(S,C),C^0(S,\mathcal O_C)\right).
\end{split}\] The profinite-product construction is an affinoid perfectoid space with underlying space \(S\); see [Retained latex: \textbackslash{}cite[Example 11.12 and Example 11.16]\{Scholze\}]. Thus \(F\) is a small \(v\)-sheaf and every component of \(|F|=\mathbf Z_2\) is a singleton. The group \(K\) is locally profinite: the trivial subgroup is an open profinite subgroup. Compactness was not required.\par
Let \(n\in\mathbf Z\) translate the \(S\)-coordinate by \(n\) and fix the coefficient-field coordinate. Pullback of continuous functions induces automorphisms of the displayed perfectoid space. This acts on the index space \(\mathbf Z_2\), not by addition in the characteristic-\(2\) coefficient field. The action is free because \(\mathbf Z\hookrightarrow\mathbf Z_2\) is injective. Its equivalence relation is \[R=\coprod_{n\in\mathbf Z}F \rightrightarrows F,
 \qquad s(n,x)=x,\quad t(n,x)=x+n.\] Both projections are pro-étale, being isomorphisms on each summand. The relation is representable and its map into \(F\times F\) is a monomorphism. Therefore \(Y=F/R=F/\underline K\) is a diamond and hence a small \(v\)-sheaf [Retained latex: \textbackslash{}cite[Definition 11.1, Proposition 11.3, Definition 12.1 and Remark 12.2]\{Scholze\}]. Quasiseparatedness is not required by that diamond definition.\par
Equation reference: eq:orbit. The quotient-topology result for a diamond presentation gives \[
 |Y|=\mathbf Z_2/\mathbf Z\] as an orbit set with its ordinary quotient topology [Retained latex: \textbackslash{}cite[Proposition 11.13 and Definition 11.14]\{Scholze\}]. This is not a Hausdorff or \(T_0\) reflection.\par
Direct calculation of the two sides. Let \(V\) be a nonempty open subset of \(\mathbf Z_2/\mathbf Z\). Its inverse image \(U\) is a nonempty open \(\mathbf Z\)-invariant subset of \(\mathbf Z_2\), so it contains \(a+2^m\mathbf Z_2\) for some \(a,m\). Integer translation runs through all residue classes modulo \(2^m\), so \[\bigcup_{n\in\mathbf Z}(a+n+2^m\mathbf Z_2)=\mathbf Z_2.\] Invariance gives \(U=\mathbf Z_2\). Thus the quotient is indiscrete and connected.\par
Equation reference: eq:contradiction. There are nevertheless at least two distinct orbits. Both \(0\) and \(1/3\) belong to \(\mathbf Z_2\), but they do not differ by an integer: \(1/3=n\in\mathbf Z\) would imply \(3n=1\) in \(\mathbf Z_2\), hence in \(\mathbf Z\) by injectivity, an impossibility. Consequently \[
 \pi_0(Y)=\{*\},\qquad
 \pi_0(F)/K=\mathbf Z_2/\mathbf Z,\qquad
 \#(\mathbf Z_2/\mathbf Z)>1.\] The claimed equality fails even as a bijection of sets. The first equality of [eq:published] (\#eq:published) holds in this example. ◻\par
\subsection*{The invalid inference}
The proof on p. 820 treats \(\pi_0\) as commuting with the relevant topological coequalizer. Applying \(\pi_0\) to the example’s action diagram leaves \(\mathbf Z\times\mathbf Z_2\rightrightarrows\mathbf Z_2\). Its coequalizer in topological spaces is the nonsingleton indiscrete space \(\mathbf Z_2/\mathbf Z\), whereas the component space of the quotient is a singleton. An adjunction into a suitable category of component objects would preserve colimits in that target category; it does not identify those colimits with ordinary orbit quotients in all topological spaces. The calculation directly isolates the failed inference.\par
\subsection*{Counterevidence and exact impact}
The quotient \(Y\) is not locally spatial: its nonsingleton indiscrete topology is not \(T_0\). The relation is not quasicompact and \(K\) is not compact. These properties do not violate the written lemma. They prevent claiming that this example occurs among the paper’s actual shtuka spaces or compact-level quotients.\par
In particular, the Iwahori quotient in formula (1.12), p. 813, and formula (6.17), p. 853, has a compact acting group. Galois-group quotient uses also involve compact groups. They are not refuted by the example. The specialization bijection equation (3.12), p. 826, belongs to Theorem 3.9 and is locally justified by rich smelted kimberlites and unibranchness. It must not be confused with Proposition 3.12 on p. 828; this review does not refute that specialization bijection.\par
The demonstrated error is a false general theorem used as the stated bridge to core component transitivity. The inspected proof supplies no replacement theorem with verified hypotheses for that bridge. Its use through Proposition 6.11 in Theorem 6.1 supports Major as a core proof-support defect. This does not assert that actual transitivity or the main theorem is false. Potential repairability does not make the demonstrated inference valid or automatically lower its severity.\par
A defensible correction must prove a component-transitivity result applicable to the particular period torsor, or give an independent argument for Proposition 3.12, and revisit affected uses. Adding local spatiality alone is not a repair established here. No universal replacement lemma is proposed.\par
\section*{A bounded unresolved check}
\subsection*{Proposition 6.6(2): converse under an ad-isomorphism}
Proposition 6.6, pp. 850–851, includes a torsor equivalence. Write \(A=G(\mathbf Q_p)\), \(B=H(\mathbf Q_p)\), \(N=\operatorname{im}G^{\mathrm{sc}}(\mathbf Q_p)\), and \(M=f(N)\). Subject to the component induction formula, the forward computation \[(A/N)\times^A B \simeq B/M\] is valid. Conversely, if a transitive \(A\)-set has stabilizer \(C\), its induced \(B\)-set has stabilizer \(f(C)\). Knowing the induced set is a \(B/M\)-torsor gives \(f(C)=M\), which alone does not give \(C=N\). Triviality of the simply connected action in the passage from (6.8) to (6.10) is immediate under the forward assumption; it does not by itself prove the converse.\par
There is substantial counterevidence to a confirmed defect. In an isotropic simply connected reduction with the required determinant control, one can obtain \(C\subset N\) and \(N=C(N\cap\ker f)\). A finite central intersection gives finite index, and the no-proper-finite-index-subgroup property can force \(C=N\). The general anisotropic situation and all geometric stabilizer hypotheses were not resolved here. No actual shtuka counterexample has been constructed.\par
This is Needs Evidence, not a second established Major or Minor finding. It asks for an explicit converse under the exact hypotheses and does not raise the supported severity or claim that every subsequently needed direction fails.\par
\section*{Checks without a separate issue}
The coarsening diagram in Section 6.3, p. 852, gives the claimed implication conditional on the stated vertical surjections and lower bijection: lift two points from the coarser level, use their common lower image, and invoke injectivity at the finer level. The set-theoretic deduction is valid.\par
The tori calculation in Lemma 6.5, p. 849, and the parahoric determinant calculation in Proposition 6.9, p. 854, are consistent with their smooth connected special-fiber and Lang-surjectivity inputs. Their formal roles were checked; external group-scheme results were not re-proved.\par
Discreteness in Lemma 6.13, p. 856, is compatible with openness of the normalizer: it contains the relevant open derived stabilizer and the center, and the central-times-derived map has open image in the \(p\)-adic group. Abbreviation of this explanation is not a separate finding. The anisotropic discussion in Theorem 6.2, p. 858, has additional geometry and is not itself disproved by the general counterexample.\par
\section*{Conclusion}
Dimension | Recorded state\par
Supported issue | Major: false Lemma 3.2 in its written generality, with an explicit core proof dependency.\par
Main theorem false? | Not established; actual geometric instances have additional structure.\par
Unresolved item | Proposition 6.6 converse: Needs Evidence.\par
Repair verification | No complete replacement transitivity proof verified.\par
Scientific adjudication | Pending; independent acceptance has not occurred.\par
Author adoption | Unknown; no author contact or adoption claim.\par
\section*{Supporting checks}
Canonical PDF SHA256:\par
> 00000TOKEN.\par
Frozen source-manifest SHA256:\par
> 00001TOKEN.\par
Reference [Scholze] is the actual MPIM PDF dated 14 April 2026, 168 pages, SHA256:\par
> 00002TOKEN.\par
It is not equated with the distinct Bonn file dated 9 September 2022, nor is byte identity with arXiv v4 asserted. The counterexample is an analytic proof, not a numerical experiment or execution of author code.\par
The primary conversation prepared and assessed this report and inspected the source evidence images. No human review is recorded.\par
These are separate attributions; independent UI observation or backend equivalence is not asserted. Actual compiler commands, dependency hashes, TeX/PDF bindings, logs and individual visual-inspection records accompany this report. Text extraction emitted a PDF syntax warning; equation-bearing source images were inspected rather than assuming every extracted glyph was accurate.\par
The coverage ledger records unperformed proof rederivations. Delivery completion is not independent acceptance.\par
9 I. Gleason, D. G. Lim and Y. Xu, The connected components of affine Deligne–Lusztig varieties, Invent. Math. 243 (2026), 805–861, online 12 November 2025. DOI https://doi.org/10.1007/s00222-025-01386-1. Publisher version of record. P. Scholze, Étale cohomology of diamonds, author-hosted manuscript dated 14 April 2026, https://people.mpim-bonn.mpg.de/scholze/EtCohDiamonds.pdf. Exact file identified above. I. Gleason, D. G. Lim and Y. Xu, same title, arXiv:2208.07195v3, 10 November 2025, https://arxiv.org/abs/2208.07195v3. Supporting source, not substituted for the publisher PDF.\par
\end{document}
